\scanpage{151}
\margin{150}
\mspage{8.}

\uline{Résumé des propositions 2 et 3.}- \uline{Le foncteur}
$\overset{\inserted{\sim}}{\operatorname{Ouv}(X)}\inserted{\xrightarrow{i_X^*}}\overset{\inserted{\sim}}{(\mathrm{Sch})}_{/X}$
\uline{est pleinement fidèle et il admet un adjoint à droite qui commute aux} $\underline{U}$-$\varinjlim$ \uline{arbitraires. De plus} $i_X^{\inserted{*}}$ \uline{commute aux} $\varprojlim$ \uline{finies.}

Tout a été démontré sauf l'assertion ``$i_X^{\inserted{*}}$ pleinement fidèle''.
Mais ceci est clair, car
\[
\operatorname{Hom}_{\overset{\inserted{\sim}}{\operatorname{Ouv}(X)}}(F,G)=\operatorname{Hom}(F,i_{X\inserted{*}}i_X^{\inserted{*}}G)=
\]
\[
=\operatorname{Hom}(i_X^{\inserted{*}}F,i_X^{\inserted{*}}G).
\]

\uline{Remarque1}.Nos raisonnements dans ce paragraphe ne sont pas particuliers à la topologie de Zariski et nous rencontrerons plus tard un autre cas où ils s'appliquent.

\uline{Remarque 2.} Dans tout ce qui précède, les faisceaux,\inserted{et les} préfaisceaux ont leurs valeurs dans la catégorie des $\underline{U}$-ensembles. Le lecteur s'assurera (en utilisant la proposition 1 ) qu'il n'y a pas de difficulté due au fait que $(\mathrm{Sch})_{/X}$ n'appartient pas à $\underline{U}$. Le foncteur $i_X^\bullet$ existe, car $\operatorname{Ouv}(X)\in\underline{U}$. On pourrait aussi faire les raisonnements pour les $\underline{V}$-(pré)-faisceaux,où $\underline{V}$ est un univers tel que $(\mathrm{Sch})_{/X}\inserted{\in}\underline{V}$, et se restreindre aux $\underline{U}$-(pré)-faisceaux après.

\S\ 2. \uline{Caractérisation ``axiomatique'' de la topologie de Zariski.}

Dans le paragraphe précédent nous avons vu que l'étude de la cohomologie de Zariski d'un préschéma $X$ peut ou bien se faire sur un grand site $(\mathrm{Sch}/X,\mathrm{Zar})$ ou bien sur des petits sites que l'on peut choisir d'une manière bien adaptée au problème en question. Dans ``la pratique'' c'est surtout ce dernier point de vue qui est dominant. Mais du point de vue ``théorique'' c'est plus agréable d'étudier la topologie de Zariski sur $(\mathrm{Sch})_{/X}$, surtout quand on veut faire des comparaisons avec d'autres topologies sur des préschémas. C'est pour préparer l'étude de ces autres topologies que nous allons donner dans ce paragraphe une caractérisation axiomatique de la topologie de Zariski ou plus précisément du petit topos de Zariski.
